\section{KZG commitments}\label{kzg}

The goal of a \alert{polynomial commitment scheme} is to have the
following API:

\begin{goal}{}{}
\begin{itemize}
\item Peggy has a secret polynomial $P(X) \in \FF_{q}[X]$.
\item Peggy sends a short ``commitment'' to the polynomial (like a hash).
\item This commitment should have the additional property that Peggy should
  be able to ``open'' the commitment at any $z \in \FF_{q}$.
  Specifically:
  \begin{itemize}
  \item Victor has an input $z \in \FF_{q}$ and wants to know
    $P(z)$.
  \item Peggy knows $P$ so she can compute $P(z)$; she sends the
    resulting number $y = P(z)$ to Victor.
  \item Peggy can then send a short ``proof'' convincing Victor that $y$
    is the correct value, without having to reveal $P$.
  \end{itemize}
\end{itemize}
\end{goal}

The \alert{Kate-Zaverucha-Goldberg (KZG)} commitment scheme\footnote{Introduced, 
of course, by Kate, Zaverucha, and Goldberg: \\ \url{https://www.iacr.org/archive/asiacrypt2010/6477178/6477178.pdf}} 
is amazingly
efficient because both the commitment and proof lengths are a single
point on $E$, encodable in 256 bits, no matter how many coefficients
the polynomial has.

\subsection{The setup}

Remember the notation $[n] \defeq n \cdot g \in E$ defined in
Definition \ref{tcb:armor}. To set up the KZG commitment scheme, a trusted party needs
to pick a secret scalar $s \in \FF_{q}$ and publish
\[ [s^0], [s^1], \dots, [s^M] \]
for some large $M$, the maximum degree of a polynomial the scheme
needs to support. This means anyone can evaluate
$[ P(s) ]$ for any given polynomial $P$ of
degree up to $M$. For example,
\[ [s^2 + 8s + 6] = [s^2] + 8[s] + 6[1].\]
Meanwhile, the secret scalar $s$ is never revealed to anyone.

The setup only needs to be done by a trusted party once for the curve
$E$. Then anyone in the world can use the resulting sequence for KZG
commitments.

\begin{remark}{}{}
The trusted party has to delete $s$ after the calculation.
If anybody knows the value of $s$, the protocol will be insecure.
The trusted party will only publish $[s^0] = [1], [s^1], \ldots, [s^M]$.
This is why we call them ``trusted'': the security of KZG
depends on them not saving the value of $s$.

Given the published values, it is (probably) extremely hard to recover the value of $s$ --
this is a case of the discrete logarithm problem.

You can make the protocol somewhat more secure by involving
several different trusted parties.
The first party chooses a random $[s_1]$,
computes $[s_1^0], \ldots, [s_1^M]$,
and then discards $s_1$.
The second party chooses $s_2$ and computes $[(s_1 s_2)^0], \ldots, [(s_1 s_2)^M]$.
And so forth.
In the end, the value of $s$ with be the product of the secrets
$s_i$ chosen by the parties... so the only way they can break secrecy
is if all the ``trusted parties'' collude.
\end{remark}

\subsection{The KZG commitment scheme}

Peggy has a polynomial $P(X) \in \FF_{p}[X]$. To
commit to it:

\begin{algo}{Creating a KZG commitment}{}
\begin{enumerate}
\item
  Peggy computes and publishes $[P(s)]$.
\end{enumerate}
\end{algo}

This computation is possible as $[s^i]$ are
globally known.

Now consider an input $z \in \FF_{p}$; Victor wants to know the
value of $P(z)$. If Peggy wishes to convince Victor that $P(z) = y$,
then:

\begin{algo}{Opening a KZG commitment}{}
\begin{enumerate}
\item
  Peggy does polynomial division to compute
  $Q(X) \in \FF_{q}[X]$ such that
  \[P(X) - y = (X - z)Q(X).\]
\item
  Peggy computes and sends Victor $[Q(s)]$,
  which again she can compute from the globally known
  $[s^i]$.
\item
  Victor verifies by checking
\begin{equation}
\label{kzg-verify}
\pair([Q(s)], [s] - [z]) = \pair([P(s)] - [y], [1])
\end{equation}
and accepts if and only if Equation \eqref{kzg-verify} is true.
\end{enumerate}
\end{algo}

If Peggy is truthful, then Equation \eqref{kzg-verify} will certainly check out.

If $y \neq P(z)$, then Peggy can't do the polynomial long division
described above. So to cheat Victor, she needs to otherwise find an
element
\[\frac{1}{s - x}( [ P(s) ] - [y] ) \in E.\]
Since $s$ is a secret nobody knows, there isn't any known way to do
this.

\subsection{Multi-openings}\label{multi-openings}

To reveal $P$ at a single value $z$, we did polynomial division to
divide $P(X)$ by $X - z$. But there's no reason we have to restrict
ourselves to linear polynomials; this would work equally well with
higher-degree polynomials, while still using only a single 256-bit curve
point for the proof.

For example, suppose Peggy wanted to prove that $P(1) = 100$,
$P(2) = 400$, \ldots, $P(9) = 8100$. (We chose these numbers so that
$P(X) = 100X^{2}$ for $X = 1,\ldots,9$.)

Evaluating a polynomial at $1,2,\ldots,9$ is essentially the same as
dividing by $(X - 1)(X - 2)\cdots(X - 9)$ and taking the remainder. In
other words, if Peggy does a polynomial long division, she will find
that
\[P(X) = Q(X)( (X - 1)(X - 2)\ldots(X - 9) ) + 100X^{2.}\]
Then Peggy sends $[Q(s)]$ as her proof, and
the verification equation is that \[\begin{aligned}
  \pair( [Q(s)], [ (s - 1)(s - 2)\ldots(s - 9) ] )  = \pair( [P(s)] - 100[s^2], [1] ).
\end{aligned}\]

The full generality just replaces the $100X^{2}$ with the polynomial
obtained from Lagrange interpolation\footurl{https://en.wikipedia.org/wiki/Lagrange\_polynomial}
(there is a unique such polynomial $f$ of degree $n - 1$). To spell
this out, suppose Peggy wishes to prove to Victor that
$P( z_{i} ) = y_{i}$ for $1 \leq i \leq n$.

\begin{algo}{Opening a KZG commitment at $n$ values}{}
\begin{enumerate}
\item
  By Lagrange interpolation, both parties agree on a polynomial $f(X)$
  such that $f( z_{i} ) = y_{i}$.
\item
  Peggy does polynomial long division to get $Q(X)$ such that
  \[P(X) - f(X) = ( X - z_{1} )( X - z_{2} )\cdots( X - z_{n} ) \cdot Q(X).\]
\item
  Peggy sends the single element $[ Q(s) ]$ as
  her proof.
\item
  Victor verifies
\[  \pair([Q(s)], [ ( s - z_1 )( s - z_2 )\ldots( s - z_n ) ]) = \pair([P(s)] - [f(s)], [1]). \]
\end{enumerate}
\end{algo}

So one can even open the polynomial $P$ at $1000$ points with a
single 256-bit proof. The verification runtime is a single pairing plus
however long it takes to compute the Lagrange interpolation $f$.

\subsection{Root check}
To make PLONK work, we're going to need a small variant of the
multi-opening protocol for KZG commitments
(Section \ref{multi-openings}), which we call
\alert{root check} (not a standard name). Here's the problem statement:

\begin{goal}{}{}
Suppose one had two polynomials $P_1$ and $P_2$,
and Peggy has given commitments
$\Com(P_1)$ and $\Com(P_2)$.
Peggy would like to prove to Victor that, say,
the equation $P_1(z) = P_2(z)$ holds for all $z$ in some large finite set $S$.
\end{goal}

Peggy just needs to show that $P_{1} - P_{2}$ is divisible by
$Z(X) \defeq \prod_{z \in S}(X - z)$. This can be done by committing the
quotient \[H(X) \defeq \frac{P_{1}(X) - P_{2}(X)}{Z(X)}.\] Victor then gives
a random challenge $\lambda \in \FF_{q}$, and then Peggy opens
$\Com( P_{1} )$,
$\Com( P_{2} )$, and $\Com(H)$ at
$\lambda$.

But we can actually do this more generally with \emph{any} polynomial
expression $F$ in place of $P_{1} - P_{2}$, as long as Peggy has a
way to prove the values of $F$ are correct. As an artificial example,
if Peggy has sent Victor $\Com( P_{1} )$ through
$\Com( P_{6} )$, and wants to show that
\[P_{1}(42) + P_{2}(42)P_{3}(42)^{4} + P_{4}(42)P_{5}(42)P_{6}(42) = 1337,\]
she could define
\[F(X) \defeq P_{1}(X) + P_{2}(X)P_{3}(X)^{4} + P_{4}(X)P_{5}(X)P_{6}(X) - 1337\]
and run the same protocol with this $F$. This means she doesn't have
to reveal any $P_{i}(42)$, which is great!

To be fully explicit, here is the algorithm:

\begin{algo}{Root check}{root-check}
Assume that $F$ is a polynomial for which Peggy can
establish the value of $F$ at any point in $\FF_{q}$. Peggy
wants to convince Victor that $F$ vanishes on a given finite set
$S \subseteq \FF_{q}$.

\begin{enumerate}
\item Peggy sends to Victor a commitment $\Com(F)$ to $F$.

\item Both parties compute the polynomial
  \[ Z(X) \defeq \prod_{z \in S}(X - z) \in \FF_{q}[X]. \]

\item Peggy does polynomial long division to compute
  $H(X) = \frac{F(X)}{Z(X)}$.

\item Peggy sends $\Com(H)$.

\item Victor picks a random challenge $\lambda \in \FF_{q}$ and
  asks Peggy to open $\Com(H)$ at $\lambda$, as well
  as the value of $F$ at $\lambda$.

\item Victor verifies $F(\lambda) = Z(\lambda)H(\lambda)$.
\end{enumerate}
\end{algo}

\begin{remark}{$\Com(F)$ may not even be necessary}{}
  In fact, rather than sending $\Com(F)$,
  it is enough for Peggy to have some way to prove to Victor the values of $F$.

  For example, if $F$ is a product of two polynomials $F = F_{1}F_{2}$,
  and Peggy has already sent commitments to $F_{1}$ and $F_{2}$,
  then there is no need for Peggy to commit to $F$.
  Instead, in Step 5, Peggy opens $\Com(F_{1})$
  and $\Com(F_{2})$ at $\lambda$, and that proves to
  Victor the value of $F(\lambda) = F_{1}(\lambda)F_{2}(\lambda)$.
\end{remark}
